% Autor: Elias Welt
\section{Softwarearchitektur und Datenfluss}
\subsection{Unidirektionaler Datenfluss}
React erzwingt einen unidirektionalen Datenfluss ("One-Way Binding"). Daten fließen immer von Eltern- zu Kind-Komponenten mittels sogenannter "`Props"'. Um Daten von einem Kind zurück zum Elternteil zu senden (z.B. ein Formular-Submit im \texttt{QuizModal}, der die Liste im \texttt{Quiz} aktualisiert), übergibt die Elternkomponente eine Callback-Funktion an das Kind.

\subsection{Routing-Strategie}
Das Routing in \texttt{App.jsx} definiert die Skelettstruktur der Anwendung. Es wird zwischen öffentlichen Routen (Login, Register) und privaten, geschützten Routen unterschieden.

\begin{lstlisting}[language=JavaScript, caption={Routing-Struktur in App.jsx}]
<Routes>
  {/* Oeffentliche Routen */}
  <Route path="/login" element={<Login setLoggedInUser={setLoggedInUser} />} />
  
  {/* Geschuetzte Routen - erfordern Login */}
  <Route path="/*" element={
    loggedInUser ? (
      <Routes>
        <Route element={<MainLayout handleLogout={handleLogout} user={loggedInUser} />}>
          <Route index element={<MainDashboard />} />
          <Route path="files" element={<Files />} />
          <Route path="chat" element={<Chat />} />
          {/* ... weitere Routen ... */}
        </Route>
      </Routes>
    ) : (
      <Navigate to="/login" />
    )
  } />
</Routes>
\end{lstlisting}

Die Verwendung von verschachtelten Routen (\texttt{<Route element={<MainLayout ...>}>}) erlaubt es, das Grundgerüst (Header, Navigation) einmalig zu definieren. Die Kind-Komponenten werden dann an der Stelle des \texttt{<Outlet />} Platzhalters in \texttt{MainLayout.jsx} gerendert.
